\chapter{Clearing Mechanism and Systemic Risk Measurement}

\section{Eisenberg--Noe Clearing Model}

Clearing is applied to payments that are contractually due on the
current business day, not to the full stock of outstanding interbank
claims. A loan with a later payment date remains part of the
outstanding-exposure matrix but contributes nothing to the current
clearing-liability matrix. This distinction allows contracts to persist
across periods without causing future principal, interest, or coupon
payments to be settled prematurely.

Suppose the banking system contains \(n\) banks. Let \(E_{ij,t}\) denote
the outstanding principal held by lender \(i\) against borrower \(j\) at
the relevant point in period \(t\). The matrix \(E_{t}\) follows a
creditor-to-debtor orientation: the first index identifies the lender
and the second identifies the borrower.

For current settlement, let \(\ A_{ij,t}^{due}\) denote the payment
scheduled from borrower \(j\) to lender \(i\). Although
\(A_{ij,t}^{due}\) follows the same orientation as the exposure matrix,
Eisenberg--Noe clearing requires liabilities to be written from debtor
to creditor. The due-payment liability matrix is therefore defined as:

\(L_{ij}^{(t)} = A_{ji,t}^{due}\).

Here, \(L_{ij}^{(t)}\) is the amount owed by debtor \(i\) to creditor
\(j\) in period \(t\). The matrix records gross contractual obligations. Reciprocal obligations
between two banks are retained as separate non-negative entries and are
not bilaterally netted before clearing. For example, if bank $i$ owes
bank $j$ while bank $j$ simultaneously owes bank $i$, both obligations
remain in $\mathbf{L}^{(t)}$ and enter the clearing calculation separately. Bank \(i\)'s total nominal payment obligation is:

\({\overline{p}}_{i,t} = \sum_{j = 1}^{n}L_{ij}^{(t)}\).

A bank with \({\overline{p}}_{i,t} = 0\) has no obligation in the current
clearing process, even if it has outstanding principal associated with
later payment dates.

The relative-liability share is defined as:


\begin{equation}
\Pi_{ij,t}=\begin{cases}\dfrac{L_{ij}^{(t)}}{\bar p_{i,t}},&\text{if }\bar p_{i,t}>0,\\[6pt]0,&\text{if }\bar p_{i,t}=0.\end{cases}\label{eq:relative-liability}.
\end{equation}


Each nonzero row of \(\Pi_{t}\) sums to one. The element
\(\Pi_{ij,t}\) represents the proportion of debtor \(i\)'s total current
payment allocated to creditor \(j\). This debtor-to-creditor orientation
differs from the lender-to-borrower orientation used for new loans and
outstanding claims in Chapter~4.

Let \(e_{i,t}\) denote the external resources available to bank
\(i\) immediately before clearing. In the baseline setting, these
resources consist of liquid assets:

\(e_{i,t} = {LA}_{i,t}^{-}\).

The implementation contains an optional \texttt{use\_core} setting under which
core capital can also be included in available resources. The baseline
equations in this chapter correspond to \texttt{use\_core=False}; any experiment
using the alternative setting should be reported separately.

Bank \(i\)'s actual payment depends on its initial liquid resources and
the payments received from other debtors. The Eisenberg--Noe clearing
condition is:

\(p_{i,t} = \min\left\{ {\overline{p}}_{i,t},e_{i,t} + \sum_{j = 1}^{n}\Pi_{ji,t}p_{j,t} \right\}\).

The incoming-payment term uses \(\Pi_{ji,t}\) because bank \(i\) is the
creditor receiving part of debtor \(j\)'s actual payment. In vector
form:

\(\mathbf{p}_{t} = \min\left\{ {\overline{\mathbf{p}}}_{t},\mathbf{e}_{t} + \Pi_{t}^{\mathsf{T}}\mathbf{p}_{t} \right\}\).

The minimum is applied component wise. No bank can pay more than its
nominal obligation, and no bank can distribute more than the resources
available from its own liquidity and incoming payments.

The clearing vector is obtained through fixed-point iteration. The
calculation begins from full nominal payment:

\(\mathbf{p}_{t}^{(0)} = {\overline{\mathbf{p}}}_{t}\).

It is then updated according to:

\(\mathbf{p}_{t}^{(k + 1)} = \min\left\{ {\overline{\mathbf{p}}}_{t},\mathbf{e}_{t} + \Pi_{t}^{\mathsf{T}}\mathbf{p}_{t}^{(k)} \right\}\).

Iteration stops when:

\(\left\| \mathbf{p}_{t}^{(k + 1)} - \mathbf{p}_{t}^{(k)} \right\|_{\infty} < \varepsilon_{EN}\).

The implementation uses a convergence tolerance of \(10^{- 6}\) and
allows at most 100 iterations.

Once the payment vector has converged, the bilateral payment from debtor
\(i\) to creditor \(j\) is:


\begin{equation}
P_{ij,t}=\Pi_{ij,t}p_{i,t}\label{eq:bilateral-payment}.
\end{equation}


The amount received by bank \(i\) is:


\begin{equation}
\mathrm{Recv}_{i,t}=\sum_{j=1}^{n}\Pi_{ji,t}p_{j,t}=\left(\boldsymbol\Pi_t^{\mathsf T}\mathbf p_t\right)_i\label{eq:received-payment}.
\end{equation}


A payment shortfall occurs when the actual payment is below the nominal
amount due:


\begin{equation}
p_{i,t}<\bar p_{i,t}-\varepsilon_{\mathrm{def}}\label{eq:clearing-default}.
\end{equation}
\vspace{-0.4\baselineskip}
\noindent where \(\varepsilon_{\mathrm{def}} = 10^{-6}\) is the numerical tolerance used to
prevent rounding errors from being classified as clearing defaults. This
condition identifies a clearing default. The distinction between
clearing default, negative equity, capital stress, and the activity
variable \(a_{i,t}\) is specified in Section~5.2.

Under the baseline resource setting, liquid assets after clearing are
updated as:


\begin{equation}
\mathrm{LA}_{i,t}^{+}=\max\left\{0,\,\mathrm{LA}_{i,t}^{-}-p_{i,t}+\mathrm{Recv}_{i,t}\right\}\label{eq:liquidity-update}.
\end{equation}


Outgoing payments reduce liquidity, while actual receipts increase it.
Creditors are therefore updated using the cash they receive through
clearing rather than the larger contractual amount originally promised.

The clearing result is then applied to the contract book. A fully paid
single-payment contract is removed after its maturity obligation has been processed.
Installment and rollover contracts retain only their updated remaining
principal and future payment schedule. The payment handled in period
\(t\) is not submitted again in a later clearing cycle.

New loans are formed only after this settlement stage. They enter the
end-of-period outstanding-exposure matrix and create payment obligations
for future dates, but they do not enter the clearing-liability matrix
for the day on which they are originated. Clearing therefore connects
the network created by earlier matching decisions with actual payment
losses while preserving the timing of each contract.

\section{Failure Definition}

In systemic risk analysis, the definition of bank failure directly
affects the calculation of the failure rate, the capital-stress share,
and the composite systemic risk indicator. If the failure definition is
too loose, the analysis may underestimate systemic risk. If the
definition is too strict, banks that have not actually defaulted but are
already under capital stress may be directly classified as failed banks,
which may overestimate the degree of default contagion. The model
distinguishes bank failure, liquidity default, and capital stress.

The simulation distinguishes an active bank, a capital-stressed bank,
and a failed bank because these states answer different questions. An
active bank may trade normally. A bank with \(\mathrm{CAR}_{i,t}<\theta\)
remains active but contributes to the capital-breach and capital-gap
measures. Hard failure occurs after negative accounting equity or, when
rollover is disabled, an unpaid single-payment obligation. Under
rollover ON, an unpaid installment becomes arrears and does not
deactivate the borrower solely because a payment was missed.

\subsection{Pre-clearing Liquidity Shortage}

Liquidity shortage is an important source of payment stress. In interbank
debt clearing, if a bank's currently available liquidity is insufficient
to cover its due payment obligations, the bank may experience a payment
shortfall. Let the available resources of bank \(i\) before clearing in
period \(t\) be \(e_{i,t}\), and its nominal payment obligation be
\({\overline{p}}_{i,t}\). If:


\begin{equation}
e_{i,t}<\bar p_{i,t}.
\end{equation}


then bank \(i\) has a liquidity shortage. In the code,
\nolinkurl{liquidity_default_candidates()} compares banks' liquid assets with
their nominal payment obligations to identify potential
liquidity-shortage banks before clearing. When \texttt{use\_core=False},
available resources mainly consist of liquid assets. When
\texttt{use\_core=True}, available resources may include both core capital and
liquid assets. In the baseline model, liquid assets are used as the payment resource for clearing.

A liquidity shortage does not, by itself, imply final failure. During
the Eisenberg--Noe clearing process, a bank may also receive payments
from other banks. Even if a bank has a shortage before clearing, it may
still complete part or all of its repayment after receiving sufficient
payments. The model therefore treats pre-clearing liquidity shortage as
a diagnostic state. Its consequences depend on the common clearing
outcome and the active rollover configuration.

\subsection{Clearing Default}

According to the Eisenberg--Noe clearing mechanism, the actual payment
of bank \(i\) is \(p_{i,t}\), and its
nominal payment obligation is \({\overline{p}}_{i,t}\). If the actual
payment is lower than the nominal obligation, the bank cannot fully meet
its interbank payment obligation:


\begin{equation}
p_{i,t}<\bar p_{i,t}.
\end{equation}


In the code implementation, if the following condition is satisfied:


\begin{equation}
p_{i,t}<\bar p_{i,t}-\varepsilon_{\mathrm{def}}.
\end{equation}


the bank has a clearing shortfall. The implementation allocates the
debtor's Eisenberg--Noe payment pro rata across all contracts due that
day. If rollover is OFF, any residual on a matured one-period contract
beyond the numerical tolerance is a hard payment default. If rollover is
ON, the residual becomes \(Q_{c,t+1}^{\mathrm{arrears}}\); consecutive
misses are recorded but do not independently deactivate the bank.
Creditors receive only actual clearing payments in either case.

\subsection{Insolvency and Negative Equity}

In addition to clearing default, a bank may also be regarded as failed
when its equity becomes negative. Negative equity represents a
balance-sheet insolvency state in which the bank's assets are
insufficient to cover its liabilities. Let the total assets of bank
\(i\) be \(TA_{i,t}\), and let its total liabilities, including current
liabilities and interbank liabilities, be \(TL_{i,t}\). Let \(E_{i,t}\) denote bank equity. It can be expressed as:


\begin{equation}
E_{i,t}=\mathrm{TA}_{i,t}-\mathrm{TL}_{i,t}.
\end{equation}


If:


\begin{equation}
E_{i,t}<0.
\end{equation}


then bank \(i\) is considered to have lost its solvency basis and can no
longer continue as a normal bank in subsequent market activities. In the
dynamic simulation process, clearing losses, investment losses, declines
in liquid assets, and increases in interbank liabilities may all worsen
bank equity. Once bank equity falls below zero, the simulation marks the
bank as inactive. The insolvency test is applied to the untruncated accounting equity
$E_{i,t}$ rather than to the stored regulatory core-capital variable.
After accounting equity has been calculated, core capital may be stored
as $\max(0,E_{i,t})$ for regulatory-ratio calculations. This non-negative
storage convention does not affect the failure test, because accounting
equity is recalculated directly from assets and liabilities when
insolvency is evaluated.

The model treats negative equity as a hard balance-sheet failure under
both rollover configurations. When a bank fails, creditor claims are
written down using the recovery rule, balance sheets and regulatory
ratios are refreshed, and newly insolvent creditors are resolved in the
same business day. This fixed-point procedure is applied after
Eisenberg--Noe settlement and again after project outcomes, so a
write-off-induced secondary failure is not postponed to the next day.

\subsection{Capital Stress}

A capital adequacy ratio below the threshold is not directly equivalent
to bank failure. In this model, the capital adequacy ratio is used to
identify whether a bank is under capital stress, rather than to directly
determine whether the bank has failed. Let the capital adequacy ratio of
bank \(i\) in period \(t\) be \(CAR_{i,t}\), and let the capital
adequacy threshold be \(\theta\). If:


\begin{equation}
\mathrm{CAR}_{i,t}<\theta.
\end{equation}


then bank \(i\) is regarded as a capital-stressed bank.

In the baseline setting, \(\theta = 0.08\). A
capital-stressed bank may still remain active and continue to
participate in the interbank market. However, its capital buffer is
already insufficient, so it is more likely to become a failed bank when
facing subsequent losses. The systemic-risk function uses inactivity to
calculate \(\mathrm{FR}_t\). CBS is calculated only among active
non-central banks. A bank that was below the CAR threshold in the
preceding period remains in CBS during a narrow recovery band until its
CAR reaches \(\theta+0.0005\).

This distinction helps describe the evolution of systemic risk in more
detail. If failure is defined only as \texttt{is\_active=False}, the recorded
results can capture only bank failures that have already occurred. If
the capital-stress share is also observed, the analysis can capture the
vulnerable part of the banking system that has not yet defaulted.
Capital stress is an early signal of systemic risk, rather than a final
failure state.

\subsection{Failure Definition in This Thesis}

Let \(F_{i,t}\) denote the absorbing failure indicator of non-central bank
\(i\), and let \(R\in\{0,1\}\) denote the rollover switch. Apart from the
small initialization failure draw reported in the experimental design,
the operative transition is:


\begin{equation}
F_{i,t}=\mathbf{1}\left\{
F_{i,t-1}=1
\ \text{or}\ E_{i,t}<0
\ \text{or}\ \bigl[R=0\ \text{and}\ U_{i,t}>\varepsilon_{\mathrm{def}}\bigr]
\right\},
\label{eq:failure-indicator}
\end{equation}

where \(U_{i,t}\) is the unpaid residual of the bank's due
single-payment contracts after clearing. When \(R=1\), unpaid installment
amounts are carried as arrears and are excluded from the third condition.

When \(F_{i,t} = 1\), the bank's activity status is set to zero:


\begin{equation}
a_{i,t}=0.
\end{equation}


The bank is then excluded from ordinary interbank matching in subsequent
periods. Obligations already recorded in the contract book remain
subject to the applicable settlement rules.

Capital stress for an active bank is defined separately. Let
\(H_{i,t}\) denote the recovery-hysteresis capital-stress indicator:


\begin{equation}
H_{i,t}=\mathbf{1}\{a_{i,t}=1\}
\begin{cases}
1, & \mathrm{CAR}_{i,t}<\theta,\\
1, & \mathrm{CAR}_{i,t-1}<\theta\ \text{and}\quad
     \mathrm{CAR}_{i,t}<\theta+0.0005,\\
0, & \text{otherwise}.
\end{cases}
\label{eq:capital-stress-indicator}
\end{equation}


A CAR breach therefore identifies an active but undercapitalized bank;
it does not by itself cause the bank to exit the market. This separation
allows \(\mathrm{FR}_t\) to measure realized failures, while \(\mathrm{CBS}_t\) and
\(\mathrm{CGR}_t\) retain information about capital weakness before failure.

\subsection{Economic Interpretation}

These conditions represent different stages of financial deterioration.
A pre-clearing liquidity shortage indicates immediate funding pressure;
a clearing shortfall is the realized unpaid amount after network receipts
have been considered; arrears preserve that amount under rollover ON;
negative equity captures balance-sheet insolvency; and a CAR breach
identifies an undercapitalized active bank. Keeping these states separate
shows whether rollover delays a liquidity failure and whether losses
ultimately create insolvency.

\section{Systemic Risk Index}

The systemic risk indicators are calculated at the end of each business
day, after scheduled obligations have been cleared, bank states have
been updated, new contracts have been registered, and the
outstanding-exposure matrix has been rebuilt. Their mathematical
definitions are given in Section~3.4.5. This section explains how the
indicators are connected to the clearing outcomes described above and
how they are used in the experimental comparison.

The composite systemic risk index is:


\begin{equation}
\mathrm{SR}_t=w_1\mathrm{FR}_t+w_2\mathrm{CBS}_t+w_3\mathrm{CGR}_t.
\end{equation}


The baseline weights are:


\begin{equation}
\left(w_1,w_2,w_3\right)=\left(0.5,0.3,0.2\right).
\end{equation}


The failure rate \(\mathrm{FR}_t\) records the proportion of non-central banks
whose activity status is zero. Under the failure rule in Section~5.2, a
non-central bank becomes inactive after negative equity or an OFF-regime
hard payment default. This component therefore represents failures that
have already materialized under the adopted classification rule.

The capital-breach share \(\mathrm{CBS}_t\) measures the share of active
non-central banks in the recovery-hysteresis low-CAR state. Failed banks
are excluded from its numerator and denominator because realized failure
is recorded by \(\mathrm{FR}_t\). The capital-gap ratio
\(\mathrm{CGR}_t\) divides aggregate capital shortfall by aggregate required
capital. For a failed bank, actual capital is set to zero and required
capital is frozen at its value at failure. Two simulations may therefore
produce similar CBS values while differing substantially in the depth of
the underlying capital deficiencies.

The three components should be examined together. An increase in
\(\mathrm{FR}_t\) indicates a rise in realized failures, whereas an increase
in \({\ CBS}_{t}\) may reveal broader vulnerability before additional
banks become inactive. Changes in \({\ CGR}_{t}\) show whether capital
shortfalls are becoming more severe, even when the number of stressed
banks changes only slightly. The composite index summarizes these
movements but does not replace the individual series.

The weights are predetermined evaluation parameters rather than values
estimated from the simulation. A larger weight on \({FR}_{t\ }\) makes
the index more responsive to realized failures, while larger weights on
\(\mathrm{CBS}_t\) or \(\mathrm{CGR}_t\) place greater emphasis on capital fragility
before or alongside default. The sensitivity analysis therefore varies
both the component weights and the CAR threshold to determine whether
the comparison between centralized matching and RFQ depends on a
particular measurement choice.

In the experimental results, \(\mathrm{SR}_t\) is reported together with
\(\mathrm{FR}_t\), \(\mathrm{CBS}_t\), and \(\mathrm{CGR}_t\). Differences between the
two matching mechanisms can then be traced to realized failures, the
breadth of capital stress, or the severity of capital shortages rather
than being inferred from the composite path alone.
